\section{Related Literature}\label{sec:literature}

The paper connects three literatures. The contribution is the interaction among them, not a claim to introduce any one component.

First, a long literature studies strategic local public inputs and mobile firms. \citet{WalzWellisch1996} analyze strategic public provision of local inputs with endogenous oligopolistic location choice. \citet{MaurerWalz2000} study two regional governments competing for two mobile oligopolistic firms and show that regional competition generally produces suboptimal infrastructure provision. Related work shows that mobility and market integration can change the direction of local-public-good distortions \citep{Matsumoto1998,Ihara2008}. Our plant-attraction motive is deliberately close to this literature. The difference is that public protection also changes firms' privately chosen geographic redundancy, so the planner's marginal valuation of the same local input becomes endogenous to the industrial resilience response.

Second, environmental and regional models make firm location and market structure endogenous to policy. \citet{MarkusenMoreyOlewiler1993} show that environmental policy can induce discrete changes in plant location and market structure, and \citet{MarkusenMoreyOlewiler1995} study noncooperative regional environmental policy when plant locations are endogenous. \citet{OtazawaEtAl2016} incorporate natural-disaster risk into an economic-geography environment and show that market equilibrium may display excessive industrial agglomeration and vulnerability. We therefore do not claim novelty for policy-induced location change or for the observation that agglomeration may create disaster exposure. Our organizational margin is instead firm-level backup readiness chosen after primary location.

Third, recent work studies private adaptation and resilience. \citet{Eisenack2014} analyzes private adaptation with endogenous market structure. \citet{CapponiDuStiglitz2024} show that market power and market incompleteness can produce insufficient supply-network resilience. \citet{CastroVincenziEtAl2024} study firms' geographic diversification of sourcing under climate risk. Most directly, \citet{ZhaoYangZhang2026} provide evidence that public flood protection can crowd out private defensive investment and increase tail-risk losses. Accordingly, public/private resilience substitution is not our novelty claim.

The paper's result is a full-game interaction. In the fixed-redundancy benchmark, the model resembles local-public-input competition. In the fixed-location benchmark, public/private resilience substitution remains but the plant-attraction motive disappears. Only the full game can generate a regime in which private redundancy adjustment makes additional common public protection socially undesirable while a local government nevertheless supplies positive protection to attract mobile oligopolistic production.
